We can say that a machine learning task is \textbf{feasible} when the following definition is satisfied.
\begin{definition}
    Let $C$ be a concept class over the instance space $X$. We say that $C$ is \textbf{PAC learnable} if and only if exists an algorithm $A$ such that
    for every $c \in C$, for every distribution $\boldsymbol{D}$, for every $0 < \epsilon < \frac{1}{2}$ and for every $0 < \delta < \frac{1}{2}$, is such that:
    $$\text{Pr}[error_{\boldsymbol{D},c}(A(EX(c, \boldsymbol{D}), \epsilon, \delta)) < \epsilon] > 1 - \delta$$
    As before, we can divide the previous formula into:
    \begin{enumerate}
        \item \textbf{Inner probability} $(error_{\boldsymbol{D},c}(A(EX(c, \boldsymbol{D}), \epsilon, \delta)))$, the probability of an error over the hypothesis $h$ is less than $\epsilon$. Therefore, $\epsilon$ is like an \textbf{upper bound} of the entire error probability.
        \item \textbf{Outer probability} $(\text{Pr}[error_{\boldsymbol{D},c}(A(EX(c, \boldsymbol{D}), \epsilon, \delta)) < \epsilon] > 1 - \delta)$, the probability that the previous error over the hypothesis $h$ is smaller than $\epsilon$ is greater than $1 - \delta$.
    \end{enumerate}
\end{definition}
\begin{definition}
    If the time complexity of $A$ is bounded by a polynomial $\frac{1}{\epsilon}$ and $\frac{1}{\delta}$, we say that $C$ is \textbf{efficiently PAC learnable}.
\end{definition}
\begin{definition}[title = Corollary]
    The concept-class of axis-aligned rectangles over $R^{2}_{[0,1]}$ is efficiently PAC-learnable.
\end{definition}